\chapter{Arbitrage de structure du capital}\label{ch:s2:capital-structure-arbitrage}

Les actions d'une entreprise et ses credit default swaps (CDS) affichent tous deux son risque de défaut. Lorsque les deux divergent, un trader achète l'un, vend l'autre et attend que le
désaccord se résorbe, ou que le modèle qui l'a mesuré se révèle faux. Yu a testé la stratégie sur 135\,759 spreads de CDS quotidiens de 261
entreprises nord-américaines, à l'aide du modèle CreditGrades. Les transactions isolées pouvaient perdre lourdement, car les spreads et les cours des actions n'étaient que faiblement corrélés. Un portefeuille
équipondéré de l'ensemble des transactions a rapporté des ratios de Sharpe comparables à ceux d'autres stratégies d'arbitrage sur titres à revenu fixe. Sur les soixante entreprises synthétiques de ce chapitre,
la stratégie affiche un ratio de Sharpe de 1.40, et 48~\% de ses transactions sont perdantes. Lorsque des rachats avec effet de levier, qui doublent la dette d'une entreprise tandis que ses actions bondissent, surviennent
quatre fois plus souvent, le ratio de Sharpe tombe à 0.39. La construction est \texttt{firm.capstruct}.

\section{Des capitaux propres au spread de crédit}

Un modèle structurel (livre~6, chapitre~14) traite les capitaux propres comme une créance sur les actifs d'une entreprise, et le défaut comme la chute des actifs sous le montant dû. La version
utilisée ici s'inspire de CreditGrades. Les actifs par action valent le cours de l'action augmenté de la part de la dette qui serait recouvrée, $V = E + LD$ pour
une dette par action $D$ et une part recouvrée $L$. La volatilité des actifs est la volatilité des capitaux propres réduite, $\sigma_E E/V$, et le défaut survient à la première date où $V$ atteint $LD$.
La probabilité de survie à cinq ans $q$, issue de la formule de premier passage, donne un spread de crédit, $(1 - R)(-\ln q)/5$, pour un taux de recouvrement du CDS $R$.

\begin{definition}[Spread implicite actions]\label{def:s2:capital-structure-arbitrage:implied}
Le \emph{spread implicite actions}\index{spread implicite actions} d'une entreprise est le spread de CDS que donne un modèle structurel à partir de son
cours de l'action, de sa volatilité des capitaux propres et de sa dette par action~; la comparaison avec le spread de marché permet de mesurer si les marchés actions et de crédit ne s'accordent pas sur
le risque de défaut de l'entreprise.
\end{definition}

Avec un cours de 50, une dette de 40 par action, une volatilité des capitaux propres de 35~\%, $L = 0.5$ et $R = 0.4$, le spread implicite à cinq ans vaut 55.8 points de base
(la \cref{fig:s2:capital-structure-arbitrage:curve}). Avec une dette de 80 par action, il vaut 115.4. Les spreads augmentent quand le cours baisse, mais seulement jusqu'à 189.8
points de base pour un cours de 10 et une dette de 40. La volatilité des actifs du modèle baisse avec les capitaux propres, une caractéristique de l'approche qui maintient les spreads en situation de détresse
à un niveau plus bas que ne le ferait un modèle à volatilité constante.

\begin{omfigure}
\begin{tikzpicture}
\begin{axis}[width=10cm,height=5.8cm,xlabel={\small cours de l'action},ylabel={\small spread à cinq ans (pb)},tick label style={font=\footnotesize},
  xmin=10,xmax=100,ymin=0,ymax=240,grid=major,grid style={gray!25},table/col sep=comma,
  legend style={font=\scriptsize,at={(0.97,0.97)},anchor=north east}]
\addplot[omThm,thick] table[x=E,y=d40]{figdata/strategies-2/09-capital-structure-arbitrage/curve.csv};
\addplot[omDef,thick,dashed] table[x=E,y=d80]{figdata/strategies-2/09-capital-structure-arbitrage/curve.csv};
\addplot[black,thick,->,>=stealth] coordinates {(50,55.8) (60,98.8)};
\legend{dette de 40 par action,dette de 80 par action}
\end{axis}
\end{tikzpicture}

\omcaption{Le \omterm{def:s2:capital-structure-arbitrage:implied}{spread implicite actions} à cinq ans du chapitre, en fonction du cours de l'action, à une volatilité des capitaux propres de 35~\%, pour une dette de 40 et de 80 par action. La flèche illustre un
rachat avec effet de levier~: le cours passe de 50 à 60 grâce à la prime, tandis que la dette double, et le spread passe de 55.8 à 98.8 points de base. Données~:
\texttt{s2\_capstruct.curve}.}\label{fig:s2:capital-structure-arbitrage:curve}
\end{omfigure}

\section{La stratégie et son ratio de couverture}

\begin{definition}[Ratio de couverture multi-actifs]\label{def:s2:capital-structure-arbitrage:hedge}
Un \emph{ratio de couverture multi-actifs}\index{ratio de couverture multi-actifs} est la quantité d'un actif qui compense la sensibilité, selon le modèle, d'une position sur
un autre actif du même émetteur~; dans l'arbitrage de structure du capital, il s'agit du nombre d'actions par unité de notionnel de CDS qui compensent la variation de valeur du CDS lorsque le
cours de l'action varie, diminué de l'annuité risquée multipliée par la dérivée du spread implicite par rapport au cours de l'action.
\end{definition}

Un trader qui constate que le spread de marché dépasse le spread implicite vend de la protection~: le CDS rémunère trop cher un risque que les capitaux propres jugent plus faible. Si le
cours de l'action baisse, le spread implicite monte et le vendeur de protection perd. Le trader vend donc aussi à découvert des actions, autant que l'indique le ratio de couverture. Pour un
cours de l'action de 50 et une dette de 40, cela représente 0.00067 action par unité de notionnel, soit environ 6\,700 actions, ou 0.34~million de dollars d'actions, face à 10~millions de dollars de
protection. Un acheteur de protection fait le contraire et achète les actions.

\texttt{firm.capstruct} applique la stratégie à soixante entreprises pendant dix ans (le \cref{lst:s2:capital-structure-arbitrage:trades}). Le spread de marché de chaque entreprise est
le spread du modèle, à la barrière et à la dette réelles de l'entreprise, multiplié par un écart de prix à retour à la moyenne (10~\% en log, demi-vie de 34 jours). Le trader ne connaît pas
la barrière réelle de chaque entreprise et suppose $L = 0.5$ partout, alors que la réalité varie d'une entreprise à l'autre avec une dispersion de 25~\%. Le trader estime la
volatilité des capitaux propres sur l'année écoulée. Il entre en position lorsque le spread de marché s'écarte de plus de 40~\% (en log) du spread implicite, et sort lorsque cet écart passe sous 5~\%
ou après 180 jours.

\section{Résultats et pertes}

\begin{center}
\small
\begin{tabular}{@{}p{5.0cm}>{\centering\arraybackslash}p{3.3cm}>{\centering\arraybackslash}p{3.3cm}>{\centering\arraybackslash}p{3.3cm}@{}}
\toprule
soixante entreprises, dix ans & rachats avec effet de levier 0.05 par entreprise-année & aucun & 0.2 par entreprise-année\\
\midrule
transactions & 440 & 413 & 433\\
moyenne par transaction (pb de notionnel) & 26.1 & 25.3 & 9.3\\
part des transactions perdantes & 48~\% & 44~\% & 52~\%\\
pire transaction (pb) & $-347$ & $-179$ & $-378$\\
part des transactions clôturées à 180 jours & 56~\% & 55~\% & 57~\%\\
ratio de Sharpe du portefeuille & 1.40 & 1.66 & 0.39\\
\bottomrule
\end{tabular}
\end{center}

Les transactions ressemblent à celles de Yu. Près de la moitié sont perdantes, les pires perdent plusieurs pour cent du notionnel, et le portefeuille qui les regroupe toutes dégage un ratio de Sharpe stable
(la \cref{fig:s2:capital-structure-arbitrage:book}). L'essentiel du rendement provient de la jambe CDS, soit 20.4 des 26.1 points de base, et la couverture en actions ajoute
5.6. Plus de la moitié des transactions ne convergent jamais et sont clôturées à 180 jours. Elles ont été ouvertes sur un écart qui n'était en réalité que la barrière erronée du trader, et non une
erreur de prix, et ces écarts ne se referment pas. Duarte, Longstaff et Yu ont constaté que les stratégies d'arbitrage sur titres à revenu fixe qui exigent davantage de «~capital intellectuel~»
ont dégagé des alphas significatifs après prise en compte des facteurs obligataires et actions et des frais, et que beaucoup présentaient des rendements à asymétrie positive, et non l'asymétrie négative de
celui qui ramasse des sous devant un rouleau compresseur.

\begin{omfigure}
\begin{tikzpicture}
\begin{axis}[width=11cm,height=5.6cm,xlabel={\small année},ylabel={\small cumulé (\% du notionnel)},tick label style={font=\footnotesize},
  xmin=1,xmax=10,ymin=-1,ymax=5,grid=major,grid style={gray!25},table/col sep=comma,
  legend style={font=\scriptsize,at={(0.03,0.97)},anchor=north west}]
\addplot[omThm,thick] table[x=year,y=base]{figdata/strategies-2/09-capital-structure-arbitrage/book.csv};
\addplot[omDef,thick,dashed] table[x=year,y=shifts]{figdata/strategies-2/09-capital-structure-arbitrage/book.csv};
\legend{rachats 0.05 par entreprise-année,rachats 0.2 par entreprise-année}
\end{axis}
\end{tikzpicture}

\omcaption{Le portefeuille synthétique d'arbitrage de structure du capital~: P\&L cumulé par transaction ouverte en moyenne (environ 27 à la fois), en part du notionnel d'une transaction,
pour des rachats avec effet de levier de 0.05 et 0.2 par entreprise-année. Données~: \texttt{s2\_capstruct.cumulative}.}\label{fig:s2:capital-structure-arbitrage:book}
\end{omfigure}

\section{Risque de modèle}

Un rachat avec effet de levier met à mal la stratégie sur ses deux jambes à la fois. La dette de l'entreprise double et son action bondit de 20~\% grâce à la prime de rachat, et le trader
apprend la nouvelle dette à partir des dépôts réglementaires trois mois plus tard. Le spread de marché bondit aussitôt. Un trader qui avait vendu de la protection perd sur le CDS, et perd
de nouveau sur les actions vendues à découvert, qui ont monté. Parmi les transactions ouvertes au moment du changement de levier de leur entreprise, les cinq qui avaient vendu de la protection ont perdu 98 points de base
en moyenne (la pire~: 316). Les sept qui avaient acheté de la protection ont gagné 220. Quand les rachats surviennent quatre fois plus souvent, 23 vendeurs de protection ont été touchés
et ont perdu en moyenne 174 points de base chacun.

Le modèle structurel est la seule mesure de l'erreur de prix sur laquelle repose la transaction, et chacune de ses entrées est incertaine~: la barrière, la volatilité des actions, la dette par action.
Un écart persistant résulte plus souvent d'une mauvaise entrée que d'une erreur de marché. Le trader a besoin de règles pour savoir quand cesser de croire au modèle. Il faut clôturer les écarts qui ne
convergent pas. Exclure les entreprises en cours de rachat. Mettre à jour la dette à partir des dépôts dès leur publication.

\section{Fiches de stratégie}

\begin{strategyfile}{Convergence du spread implicite actions}\label{strat:s2:capital-structure-arbitrage:convergence}
\sfield{Qui paie, et pourquoi} Des investisseurs d'un marché qui valorisent le risque de défaut d'une entreprise autrement que les investisseurs de l'autre marché, jusqu'à ce que les deux
prix convergent.
\sfield{Instruments et plateformes} CDS sur un nom unique~; les actions de l'entreprise pour la couverture.
\sfield{Signal} Le spread de marché comparé au \omterm{def:s2:capital-structure-arbitrage:implied}{spread implicite actions} issu d'un modèle structurel.
\sfield{Taille et exécution} De nombreuses petites positions~; ratio de couverture issu du modèle~; sortie à la convergence ou après une limite de temps.
\sfield{Coûts} Fourchettes des CDS~; emprunt des actions vendues à découvert.
\sfield{Comment elle meurt} Les écarts qui relèvent d'une erreur de modèle~; les rachats et les recapitalisations qui font bouger l'action et le crédit dans le même sens.
\sfield{Horizon, capacité, infrastructure} Quelques mois~; un modèle structurel alimenté par les dépôts, la volatilité des actions et les cotations de CDS.
\sfield{Backtest honnête} Dette telle qu'elle était connue à l'époque~; inclure les transactions clôturées à la limite de temps.
\sfield{Sources} Yu (2006)~; ce chapitre~: un ratio de Sharpe du livre de 1.40 et 48~\% des transactions perdantes.
\end{strategyfile}

\begin{strategyfile}{Achat de CDS, vente à découvert d'actions faibles}\label{strat:s2:capital-structure-arbitrage:weak}
\sfield{Qui paie, et pourquoi} Des investisseurs en crédit plus lents que les investisseurs en actions à intégrer les mauvaises nouvelles.
\sfield{Instruments et plateformes} Protection CDS achetée~; actions vendues à découvert.
\sfield{Signal} Un cours de l'action en baisse avec un spread de CDS inférieur à son niveau implicite.
\sfield{Taille et exécution} Dimensionnée selon le ratio de couverture~; petite sur les noms proches de la détresse.
\sfield{Coûts} Prime de protection~; emprunt des actions difficiles à emprunter.
\sfield{Comment elle meurt} Reprise du cours des actions~; rappel de l'emprunt.
\sfield{Horizon, capacité, infrastructure} De quelques semaines à quelques mois.
\sfield{Backtest honnête} Coûts d'emprunt sur les actions en baisse.
\sfield{Sources} Aucune performance vérifiée.
\end{strategyfile}

\begin{strategyfile}{Désalignement dette--actions après des événements}\label{strat:s2:capital-structure-arbitrage:events}
\sfield{Qui paie, et pourquoi} Des détenteurs contraints de vendre un côté après un événement (une dégradation, une exclusion d'indice, l'annonce d'un rachat).
\sfield{Instruments et plateformes} CDS, obligations et actions du même émetteur.
\sfield{Signal} L'événement et le sens dans lequel il fait bouger chaque créance.
\sfield{Taille et exécution} Petite~; l'événement peut modifier la structure du capital elle-même.
\sfield{Coûts} Fourchettes élargies autour des événements.
\sfield{Comment elle meurt} L'événement est un changement de levier~: l'action et le crédit sont tous deux réévalués, correctement.
\sfield{Horizon, capacité, infrastructure} De quelques jours à quelques semaines~; suivi des événements.
\sfield{Backtest honnête} Dates des événements et dépôts tels qu'ils étaient connus à l'époque.
\sfield{Sources} Aucune performance vérifiée~; ce chapitre~: les vendeurs de protection ouverts lors d'un rachat ont perdu 98 points de base en moyenne.
\end{strategyfile}

\section{Tutoriel~: même défaut, deux prix}\label{tut:s2:capital-structure-arbitrage:tut}

\begin{tutorial}
\textbf{Objectif.} Simuler des entreprises dont le CDS et l'action sont régis par un modèle structurel, avec du bruit introduit volontairement et des changements de levier, négocier la convergence et mesurer
les pertes lorsque le modèle se trompe. \textbf{État final~:} le tableau et les deux figures.
\begin{tutsteps}
\item \textbf{Les transactions}.
\omcode{code/firm/capstruct/firm_capstruct.py}{118}{154}{Entrée sur l'écart, jambes CDS et actions jour par jour, sortie à la convergence ou à l'échéance.}\label{lst:s2:capital-structure-arbitrage:trades}
\item \textbf{Résultats}.
\omcode{code/strategies-2/09-capital-structure-arbitrage/python/s2_capstruct.py}{33}{54}{Statistiques par transaction et du livre, et transactions touchées par un changement de levier.}\label{lst:s2:capital-structure-arbitrage:stats}
\item \textbf{Exécuter} \texttt{stats()}, \texttt{stats(0.2)}, \texttt{through\_shifts()}, \texttt{curve()} et \texttt{fig\_capstruct.py}.
\end{tutsteps}
\textbf{Pour aller plus loin.} Donner au trader la vraie barrière et mesurer la part de la perte due à l'erreur de modèle~; écarter les entreprises dont l'action a bondi de plus de
15~\% en une journée~; n'entrer qu'après un écart persistant pendant un mois.
\end{tutorial}

\section{Construction~: arbitrage de structure du capital}\label{bld:s2:capital-structure-arbitrage:capstruct}

\begin{build}
\bfield{Rôle} \omterm{def:s2:capital-structure-arbitrage:implied}{Spreads implicites actions} issus d'un modèle de premier passage, ratios de couverture, transactions de convergence et test de stress d'un changement de levier.
\bfield{Interface} \texttt{CapConfig(\dots)}, \texttt{implied\_spread(E, D, sigma\_E, cfg)}, \texttt{hedge\_ratio(E, D, sigma\_E, cfg)},
\texttt{simulate\_firms(cfg)}, \texttt{trades(sim, cfg)}.
\bfield{Règles} Le trader voit la dette avec un décalage et suppose une barrière unique pour toutes les entreprises~; un notionnel unitaire par transaction~; coûts à l'entrée et à la sortie.
\bfield{Tests d'acceptation} \texttt{code/firm/capstruct/tests/}~: les spreads montent avec le levier et baissent avec l'action~; la formule de survie vérifiée à la main~; sans
erreur de prix ni erreur de modèle, aucune transaction n'est ouverte.
\bfield{Pour aller plus loin} Incertitude sur la barrière comme dans CreditGrades~; jambes obligataires~; risque de recouvrement.
\end{build}

\begin{omsources}
\item F.~Yu, ``How profitable is capital structure arbitrage?'', \emph{Financial Analysts Journal} 62(5), 2006.
\item J.~Duarte, F.~A.~Longstaff et F.~Yu, ``Risk and return in fixed-income arbitrage: nickels in front of a steamroller?'', \emph{Review of
Financial Studies} 20(3), 2007.
\end{omsources}

\section{Exercices}

\begin{exercise}[$\star$]\label{exo:s2:capital-structure-arbitrage:1}
Une probabilité de survie à cinq ans est de 0.954 et le taux de recouvrement est de 40~\%. Quel spread la formule du chapitre donne-t-elle~?
\end{exercise}

\begin{exercise}[$\star$]\label{exo:s2:capital-structure-arbitrage:2}
Le ratio de couverture est de 0.00067 action par unité de notionnel et le cours de l'action est de 50. Quelle quantité d'actions couvre 10~millions de dollars de protection~?
\end{exercise}

\begin{exercise}[$\star$]\label{exo:s2:capital-structure-arbitrage:3}
Un vendeur de protection avec une annuité risquée de 4.5 voit le spread passer de 56 à 99 points de base. Combien perd-il par unité de notionnel~?
\end{exercise}

\begin{exercise}[$\star\star$]\label{exo:s2:capital-structure-arbitrage:4}
Pourquoi plus de la moitié des transactions synthétiques se clôturent-elles à la limite de temps plutôt qu'à la convergence~?
\end{exercise}

\begin{exercise}[$\star\star$]\label{exo:s2:capital-structure-arbitrage:5}
Pourquoi un rachat avec effet de levier nuit-il au vendeur de protection sur ses deux jambes~?
\end{exercise}

\begin{exercise}[$\star\star$]\label{exo:s2:capital-structure-arbitrage:6}
Yu a trouvé de lourdes pertes sur certaines transactions isolées et un portefeuille régulier. Comment les deux peuvent-ils être vrais~?
\end{exercise}

\begin{exercise}[$\star\star\star$]\label{exo:s2:capital-structure-arbitrage:7}
\emph{Programmation.} Exécuter \texttt{stats(0.0)}. Que fait la suppression des changements de levier à la pire transaction et au ratio de Sharpe, et quel risque subsiste~?
\end{exercise}

\begin{exercise}[$\star\star\star$]\label{exo:s2:capital-structure-arbitrage:8}
\emph{Trouver la faille.} «~Le CDS est supérieur de 50~\% à son niveau implicite actions depuis six mois, donc la transaction de convergence est plus sûre que jamais.~»
\end{exercise}

\section{Problème~: Même défaut, deux prix}

\begin{problem}[{Problème du week-end --- arbitrage de structure du capital}]\label{pb:s2:capital-structure-arbitrage:1}
Les entreprises synthétiques du chapitre et les données publiques.

\medskip\noindent\textbf{Partie I --- Le modèle.}
\begin{enumerate}
\item Décrire le modèle de premier passage et ses données d'entrée.
\item Définir le \omterm{def:s2:capital-structure-arbitrage:implied}{spread implicite actions} et le calculer pour un cours d'action de 50, avec une dette de 40 et de 80.
\item Pourquoi les spreads du modèle restent-ils modérés lorsque l'action baisse~?
\item Que Yu a-t-il utilisé et trouvé~?
\end{enumerate}

\medskip\noindent\textbf{Partie II --- La transaction.}
\begin{enumerate}[resume]
\item Définir le \omterm{def:s2:capital-structure-arbitrage:hedge}{ratio de couverture multi-actifs} et calculer la couverture pour 10~millions de dollars.
\item Décrire les règles d'entrée et de sortie.
\item Qu'ignore le trader sur chacune des entreprises~?
\item D'où provient l'essentiel du rendement~?
\end{enumerate}

\medskip\noindent\textbf{Partie III --- Les résultats.}
\begin{enumerate}[resume]
\item Donner les statistiques du portefeuille avec des rachats d'entreprises à 0.05 par entreprise-année.
\item Pourquoi tant de transactions arrivent-elles à leur date limite~?
\item Qu'ont établi Duarte, Longstaff et Yu~?
\item Comment le portefeuille se compare-t-il à une situation sans rachats~?
\end{enumerate}

\medskip\noindent\textbf{Partie IV --- Le verdict.}
\begin{enumerate}[resume]
\item Énoncer le \emph{résultat nommé}~: le ratio de Sharpe de la transaction de convergence et les pertes lorsque l'effet de levier varie.
\item Qu'est-il advenu des transactions ouvertes lors d'un rachat, selon leur côté~?
\item Pourquoi la barrière de défaut est-elle la donnée d'entrée la plus faible du modèle~?
\item Quelles règles faudrait-il ajouter~?
\item Comment mener honnêtement un backtest de la transaction~?
\item Quelle fiche de stratégie est le plus exposé aux rachats~?
\item Quel lien ce chapitre entretient-il avec le modèle structurel du livre~6~?
\item En une phrase~: sur quoi parie un arbitragiste de structure du capital~?
\end{enumerate}
\end{problem}

\section{Questions d'entretien}

\begin{interviewq}[$\star$ \iqroles{chercheur}]\label{iq:s2:capital-structure-arbitrage:1}
Comment obtiendriez-vous un spread de crédit à partir du cours d'une action~?
\end{interviewq}

\begin{interviewq}[$\star\star$ \iqroles{trader}]\label{iq:s2:capital-structure-arbitrage:2}
Le CDS d'une entreprise est nettement supérieur à son niveau implicite actions. Décrivez la transaction et ce qui pourrait mal tourner.
\end{interviewq}

\begin{interviewq}[$\star\star$ \iqroles{risque}]\label{iq:s2:capital-structure-arbitrage:3}
Comment soumettriez-vous un portefeuille d'arbitrage de structure du capital à des tests de résistance~?
\end{interviewq}

\begin{interviewq}[$\star\star$ \iqroles{chercheur}]\label{iq:s2:capital-structure-arbitrage:4}
Comment distingueriez-vous une mauvaise valorisation d'une erreur de modèle~?
\end{interviewq}

\begin{interviewq}[$\star\star$ \iqroles{développeur}]\label{iq:s2:capital-structure-arbitrage:5}
Quelles données le portefeuille exige-t-il chaque jour, et lesquelles arrivent en retard~?
\end{interviewq}

\begin{interviewq}[$\star\star\star$ \iqroles{chercheur}]\label{iq:s2:capital-structure-arbitrage:6}
Pour des actifs suivant un mouvement brownien géométrique sans dérive (en termes logarithmiques, de dérive $-\sigma^2/2$) et une barrière constante $B$, établir la probabilité de
ne pas toucher la barrière jusqu'à l'instant $t$, en utilisant le principe de réflexion.
\end{interviewq}
